% Constructive atomless branch; no claim of original-edition admission.

\begin{annexitem}{annex:p17-sa-cup-consistency}{thm:p17-sa-cup-consistency}{Qualitative additivity, copying and consistent CUP cuts}
\begin{proof}
Fix stakes \(M\succ m\). P4 makes event comparisons independent
of this choice. P1 makes them a weak order. An event is equivalent
to the empty event exactly when it is null: a nonnull event,
by P3, strictly prefers its \(M\)-over-\(m\) replacement.
Inclusion is weakly monotone; addition of a nonnull difference
is strict. P2 cancels a common winning event or inserts it
where both bets previously lost, preserving weak, strict and
indifferent comparisons. Its constant-background version suffices
on simple acts, by changing each cell of a finite common-background
partition successively; the reverse implication is immediate.

We spell out the source's qualitative-additivity diagram.
Suppose \(A_1\succeq_e B_1\), \(A_2\succeq_e B_2\), with
each pair of families disjoint internally. Put
\(A_0=(A_1\cup A_2)^c\), \(B_0=(B_1\cup B_2)^c\),
and \(C_{ij}=A_i\cap B_j\).
After cancelling common parts the premises are
\[
 C_{10}\cup C_{12}\succeq_e C_{01}\cup C_{21},\qquad
 C_{20}\cup C_{21}\succeq_e C_{02}\cup C_{12}.
\]
Add common \(C_{20}\) to the first and \(C_{01}\) to the
second. Their middle events coincide, so transitivity gives
\[
 C_{10}\cup C_{20}\cup C_{12}
 \succeq_e C_{01}\cup C_{02}\cup C_{12}.
\]
Cancel \(C_{12}\), then restore the common intersection of
the two original unions. This proves
\(A_1\cup A_2\succeq_e B_1\cup B_2\), strictly if a
premise was strict. Iteration proves finite additivity of
comparisons and equalities. If equal total events have equal
subevents, their differences are equal: otherwise strict
additivity contradicts equality of the totals. In particular
equal events have equal complements.

If \(C\succeq_e A\succ_e\varnothing\), there is
\(D\subseteq C\) with \(D\sim_e A\).
At equality use \(D=C\). Otherwise the bet on \(A\) lies
strictly between constant \(m\) and the bet on \(C\), so
P6.2 supplies exactly such a copy.
Starting with \(A\), repeatedly copy it into the remainder
while that remainder is at least as likely. P6.1 precludes
an infinite disjoint equally likely nonnull sequence,
so this produces a finite CUP.
No CUP of \(n\) units can accommodate \(n+1\) disjoint
copies of its unit anywhere in \(S\): the first \(n\) copies
are equivalent to the CUP's \(n\) units, hence their
complements are equivalent; the new complement contains
a unit and cannot be less likely than it, whereas the
original remainder is less likely.

There are arbitrarily fine CUPs. Split a unit \(A\) into
two nonnull parts using absence of atoms; choose the less
likely part \(B\). The other part is at least as likely
as \(B\). In each unit \(A_j\) of a given CUP, copying
and cancellation give two disjoint copies of \(B\).
Extend these \(2n\) copies greedily to a CUP with at least
\(2n\) units. Repetition gives arbitrarily many.
The recursive existential choices use the usual classical
choice framework, not an effective event-selection oracle.

Fix a coarse CUP of \(n\) units and weak cut \(k\) for \(E\).
For a fine CUP of \(N\) units \(B_j\), let \(i\) be the
largest number whose initial union is no more likely than
the coarse unit \(A_1\), and \(\ell\) the corresponding
weak cut for \(E\).
Copying \(i\) fine units into each coarse unit shows
\(ni\leq N\), by maximality of a fine CUP.
If \(N\geq(n+1)(i+1)\), arrange \(n+1\) disjoint groups
of \(i+1\) original fine units. Each exceeds a coarse unit,
and also exceeds its remainder. Strict qualitative additivity
would make their union exceed \(S\), although it is a subset
of \(S\). Hence
\[
 ni\leq N<(n+1)(i+1).
\]
If \(i=N\), the upper inequality is immediate, so it does
not require a nonexistent next fine cell.

Similarly copy \(i\) fine units into each of the \(k\)
coarse units below \(E\), obtaining \(\ell\geq ki\).
If \(k<n\), \(E\) is below the next coarse cumulative
union. Groups of \(i+1\) fine units, when available, exceed
its constituent coarse units, giving
\(\ell<(k+1)(i+1)\). If there are insufficient fine units,
the same weak bound follows from \(\ell\leq N\).
For \(k=n\), it follows from \(E\subseteq S\) and the
already proved upper bound on \(N\). Thus
\[
 ki\leq\ell\leq(k+1)(i+1).
\]
As \(N\to\infty\), \(i\to\infty\), since otherwise the
upper bound on \(N\) would remain bounded.
For positive \(i\) the conservative fine bounds satisfy
\[
 \frac{\ell-1}{N+1}
 \geq\frac{ki-1}{(n+1)(i+1)}
 \longrightarrow\frac{k}{n+1}>
 \frac{k-1}{n+1},
\]
\[
 \frac{\ell+2}{N}
 \leq\frac{(k+1)(i+1)+2}{ni}
 \longrightarrow\frac{k+1}{n}<
 \frac{k+2}{n}.
\]
This includes \(k=0\); the first lower bound tends to zero.
Therefore every sufficiently fine interval lies inside
the fixed coarse interval. For any finite collection of
coarse intervals choose a fine CUP large enough for all
of them simultaneously. Their intersection is nonempty.
The finite-intersection property inside one fixed compact
interval gives a point common to all CUP intervals.
Their widths are at most \((4n+2)/(n(n+1))\), tending to
zero, so that point is unique.
The numbers \(k/(n+1)\) lie in those intervals and in
\([0,1]\); their fine-CUP limit is \(P(E)\).
The empty event has cut zero and the whole space cut \(n\),
giving \(P(\varnothing)=0\), \(P(S)=1\).

For additivity, a CUP unit can be copied greedily into any
event \(E\) exactly \(k_E\) times, leaving a remainder less
likely than a unit. This follows by cancellation at each
step from the definition of its maximal weak cut.
For disjoint \(E,F\), the two sets of copies are disjoint.
Their two remainders together are less likely than two
units, by qualitative additivity, and so can supply at most
one further unit. For CUPs with at least two units,
\[
 k_E+k_F\leq k_{E\cup F}\leq k_E+k_F+1.
\]
Divide by \(n+1\) and take the fine-CUP limit:
\(P(E\cup F)=P(E)+P(F)\).

Event order implies order of every cut, hence weak order
of probabilities. A nonnull \(D\) generates a CUP with
\(m\) copies of \(D\) and a smaller remainder. Weak order
preservation and finite additivity give
\(1=mP(D)+P(R)\leq(m+1)P(D)\), so \(P(D)>0\).
If \(A\succ_e B\), copy \(B\) into \(A\) (use the empty
event if \(B\) is null). The difference is nonnull,
or qualitative additivity would imply equality.
Consequently \(P(A)>P(B)\).
Completeness gives the reverse representation direction.

Finally any normalized finitely additive probability \(Q\)
representing the event order assigns a CUP unit mass
\(a\in[1/(n+1),1/n]\), since the remainder is no larger.
Its cut sandwiches \(Q(E)\) between \(ka\) and
\((k+1)a\), with \(Q(E)\leq1\) in the last-cell case.
Thus \(Q(E)\) belongs to every conservative CUP interval.
Their unique intersection gives \(Q(E)=P(E)\).
\end{proof}
\end{annexitem}

\begin{annexitem}{annex:p17-sa-utility}{thm:p17-sa-utility}{Conditional exchange rates and the full simple-act construction}
\begin{proof}
For \(m\prec x\prec M\), P3 and solvability provide a
conditional standard gamble on every nonnull \(A\);
at tied endpoints choose the empty or full event.
For a null event choose the empty event.
P2 makes calibration independent of the outside background.
Two calibrations on the same event yield indifferent
\(M/m\) bets, hence equally likely winning events by P4.
Therefore \(T_x(A)=P(A^x)\) is well defined.
It has \(0\leq T_x(A)\leq P(A)\).

For disjoint \(A,B\), replace the \(x\) consequence on
each by its conditional standard gamble, one at a time,
using P2. The resulting winning events are disjoint.
Comparing with the calibration on \(A\cup B\) gives
\(T_x(A\cup B)=T_x(A)+T_x(B)\).
Equal-likelihood events have equal \(T_x\) values: compare
the \(x/m\) bets by P4, then replace them by their
calibrated \(M/m\) bets. Copying a smaller event into a
larger one and using positivity shows that \(T_x\) also
preserves weak event order.

Put \(u=T_x(S)\). On a CUP let
\(a=P(A_1)\), \(t=T_x(A_1)\), \(r=P(R)\),
\(t_R=T_x(R)\). Then
\[
 1=na+r,\quad u=nt+t_R,\quad
 0\leq t\leq a\leq1/n,\quad 0\leq t_R\leq r<a.
\]
It follows that
\[
 \left|\frac ta-u\right|
 =\left|\frac{ur-t_R}{1-r}\right|
 \leq\frac r{1-r}\longrightarrow0.
\]
For an arbitrary event \(E\) with CUP cut \(k\),
event-order preservation sandwiches both \(T_x(E)\) and
\((t/a)P(E)\) between \(kt\) and \((k+1)t\).
For the last-cell case use \(T_x(R)\leq t\).
Their difference is at most \(t\leq a\to0\).
Therefore \(T_x(E)=uP(E)\) for every event.
Taking \(U(x)=T_x(S)\) proves the conditional calibration
identity, including the endpoints.

In any finite-valued act between \(m,M\), replace each
consequence on its event by its conditional standard gamble.
P2 preserves each indifference and P1 composes them.
The resulting act wins \(M\) on the disjoint union of its
calibrating subevents and \(m\) elsewhere, with probability
\(\sum_jP(E_j)U(x_j)\).
Comparing two such binary acts by the represented event
order proves both directions of the SEU formula.
Constant comparisons show that \(U\) preserves consequence
order and is nonconstant.

For any finite collection of consequences, choose its
least and greatest elements under the complete consequence
order, including the fixed \(m,M\). Repeat the construction
on that larger order interval and normalize its utility
to zero at \(m\), one at \(M\).
The normalized representations agree on every overlap.
Indeed for \(m\preceq x\preceq M\) they all give
\(U(x)=P(S^x)\).
For \(x\succ M\), solvability gives
\(M\sim x_Em\) with \(P(E)>0\), forcing
\(U(x)=1/P(E)\) in every normalized representation.
For \(x\prec m\), it gives
\(m\sim M_Ex\) with \(P(E)<1\), forcing
\(U(x)=-P(E)/(1-P(E))\).
Ties have utility zero or one as appropriate.
Hence these locally constructed utilities patch to one
function on all \(X\), and any pair of simple acts lies
in one finite collection covered by this construction.
Every other representing probability is already the same
by the preceding annex proof. After fixing the two anchors,
the same calibration equations force every other utility
to agree. Undoing normalization gives precisely positive
affine freedom. No global boundedness was inferred.
\end{proof}
\end{annexitem}

\begin{annexitem}{annex:p17-sa-general-converse}{thm:p17-sa-general-converse}{Bounded finitely additive integration and P1--P7}
\begin{proof}
Refinement of finite partitions makes the integral of a simple
function independent of its description, by finite additivity.
Common refinement proves linearity and positivity.
It also gives the norm bound
\(|I_P(v)-I_P(w)|\leq\lVert v-w\rVert_\infty\) for simple
functions. Every bounded real function has finite-valued
uniform approximations, for instance rounding down to a
dyadic grid. Their integrals form a Cauchy sequence.
The norm bound proves independence of the approximation
sequence and gives a unique continuous extension.
Linearity, positivity, monotonicity and the norm bound
pass to that extension. Localization gives
\[
 |I_P(\mathbf1_A(v-w))|
 \leq P(A)\lVert v-w\rVert_\infty.
\]
This uses no countable-additivity premise.

Compare all acts by \(V(f)=I_P(u\circ f)\), with
bounded nonconstant \(u\) and convex-ranged \(P\).
The numerical order gives P1.
P2 follows because the common outside contribution cancels.
An event is null exactly when \(P(A)=0\):
zero mass makes every localized difference zero; positive
mass and two differently valued consequences give a strict
conditional comparison. On a nonnull event the constant
difference is \(P(A)(u(x)-u(y))\), proving P3.
Two-prize event bets have difference
\((u(x)-u(y))(P(A)-P(B))\), proving P4.
Nonconstancy gives P5.

Let the utility range lie in \([a,b]\), \(b>a\), and
\(V(f)-V(g)=d>0\).
Convex-rangedness constructs a partition with \(N\) equal
probability cells by repeated extraction of mass \(1/N\).
Choose \(N\) with \((b-a)/N<d\).
Replacing either act on one cell by any constant consequence
changes its value by at most that bound, proving P6.
The same proof even permits an arbitrary bounded-utility
replacement act; this stronger conclusion does not redefine P6.

For P7 on a nonnull \(A\), set
\(c_f=I_P(\mathbf1_Au(f))/P(A)\), and similarly \(c_g\).
If every constant \(f(s)\succeq_Ag\), then
\(u(f(s))\geq c_g\) on \(A\).
Positivity of the integral gives \(c_f\geq c_g\).
If \(f\succeq_Ag(s)\) for every \(s\in A\), then
\(c_f\geq u(g(s))\) on \(A\), giving the same conclusion.
On a null event both claims are automatic.
This proves the two actual P7 clauses for all acts.
\end{proof}
\end{annexitem}

\begin{annexitem}{annex:p17-sa-identification}{thm:p17-sa-identification}{Identification and explicit failures}
\begin{proof}
Suppose \((P,u)\), \((Q,v)\) are state-independent
representations of the same full act order with convex-ranged
probabilities and nonconstant utilities.
Common-stake bets induce the same event order.
A \(P\)-uniform partition into \(n\) cells has all its
cells equally likely; the same is true under \(Q\).
Normalization makes each \(Q\) cell have mass \(1/n\).
Their unions realize every rational \(k/n\) under both.
Comparing an arbitrary event with these unions forces
the same rational cuts for \(P(E)\) and \(Q(E)\);
a rational strictly between two unequal real numbers would
contradict that comparison. Thus \(P=Q\).

Normalize both utilities at \(M\succ m\).
For \(0\leq u(x)\leq1\), convex-rangedness gives a bet
winning on mass \(u(x)\) and indifferent to constant \(x\),
so \(v(x)=u(x)\).
For \(u(x)>1\), use a winning mass \(1/u(x)\) and
\(x_Em\sim M\), giving the same equality.
For \(u(x)<0\), use losing mass \(1/(1-u(x))\) and
\(x_EM\sim m\). Its expected value equation again gives
\(v(x)=u(x)\). Undo normalization to obtain affine uniqueness.

A probability atom \(A\) of positive mass cannot meet P6.
Compare \(M_Am\succ m\) and replace by \(m\).
In any finite partition, one cell contains all of \(A\)'s
probability mass: every intersection has mass zero or
\(P(A)\), and their masses sum to \(P(A)\).
The replacement on that cell leaves only a null winning
event and is indifferent to \(m\), contradicting P6.

In the two-state binary-consequence example both probability
vectors rank the losing constant, the second-state win,
the first-state win and the winning constant in the same
strict order, despite their different numbers.
For the state-dependent example let
\(r=1.5\mathbf1_E+.5\mathbf1_{E^c}\).
Finite additivity defines \(Q(A)=I_P(r\mathbf1_A)\),
with \(Q(S)=1\) and \(Q(E)=3/4\).
The bounded function \(u(f)/r\) integrates under \(Q\)
to \(I_P(u(f))\), first for simple functions by splitting
on \(E,E^c\), then for all bounded functions by uniform
extension. Thus altering the state-dependent scale can
offset the altered probability. It is outside the
state-independent uniqueness theorem, not its exception.
\end{proof}
\end{annexitem}
